\subsection{Analisis Kebutuhan Sistem}
\label{subsec:analisis-kebutuhan-sistem}
Untuk mengatasi masalah pada subbab \ref{subsec:identifikasi-masalah}, sistem memerlukan serangkaian kebutuhan fungsional yang selaras dengan rumusan masalah, tujuan penelitian, dan skenario layanan kesehatan kolaboratif pada Subsubbab \ref{subsubsec:analisis-skenario-layanan-klinis}. Secara khusus, kebutuhan sistem harus mampu mengakomodasi perbedaan peran antara petugas administrasi, dokter penanggung jawab pasien, dan petugas laboratorium, yang masing-masing membutuhkan ruang lingkup akses berbeda terhadap data RME. Kebutuhan tersebut dirangkum dalam Tabel \ref{tab:kebutuhan-fungsional-tambahan}.

\begin{table}[!ht]
	\centering
	\caption{Pemetaan Masalah dengan Kebutuhan Fungsi Tambahan}
	\label{tab:kebutuhan-fungsional-tambahan}
	\small
	\begin{tabular}{|c|p{0.18\textwidth}|p{0.2\textwidth}|p{0.38\textwidth}|}
		\hline
		\textbf{Masalah} & \textbf{Fungsionalitas}                            & \textbf{Aktor dan Kebutuhan Akses}                                                                                                                                                        & \textbf{Deskripsi Kebutuhan}                                                                                                                                         \\ \hline
		1                & Kontrol Akses Granular Berbasis Macaroons          & Petugas administrasi hanya memerlukan akses administratif; dokter memerlukan akses baca/tulis klinis; petugas laboratorium hanya memerlukan akses terbatas sesuai permintaan pemeriksaan. & Sistem harus mampu menerbitkan token otorisasi yang membawa batasan rinci berdasarkan aktor, jenis data, aksi, masa berlaku, dan konteks penggunaan akses.           \\ \hline
		2                & Segmentasi Data \textit{Off-Chain} Terenkripsi     & Data administratif, riwayat klinis, permintaan pemeriksaan, dan hasil laboratorium perlu dapat diakses secara selektif oleh aktor yang berbeda.                                           & Sistem harus memecah data RME ke dalam unit penyimpanan \textit{off-chain} yang tersegmentasi agar setiap aktor hanya dapat mengakses himpunan data yang relevan.    \\ \hline
		3                & Delegasi Berantai dengan atenuasi \textit{offline} & Dokter perlu dapat mendelegasikan sebagian hak akses kepada petugas laboratorium tanpa meminta pasien menerbitkan ulang token.                                                            & Sistem harus memungkinkan tenaga medis mendelegasikan sebagian hak akses secara lokal dengan batasan yang makin ketat tanpa memerlukan intervensi ulang dari pasien. \\ \hline
	\end{tabular}
\end{table}

Dengan demikian, keluaran tahap analisis pada bab ini menjadi dasar langsung bagi tahap perancangan solusi, yaitu integrasi Macaroons sebagai mekanisme kontrol akses yang dapat membedakan kebutuhan akses setiap aktor, penyusunan skema penyimpanan \textit{off-chain} yang tersegmentasi sesuai unit data klinis, dan perancangan alur delegasi berantai yang mendukung perpindahan hak akses secara efisien dalam layanan kesehatan kolaboratif.
